% Agent Behavior Regression Testing for Local LLM Systems
% Seed Reproducibility and Cross-Model Heterogeneity across 100 Scenarios
% TMLR Submission - official JmlrOrg/tmlr-style-file template
\documentclass[11pt]{article}
\usepackage[preprint]{tmlr}
\input{math\_commands}

\usepackage[T1]{fontenc}
\usepackage[utf8]{inputenc}
\usepackage{graphicx}
\usepackage{amsmath,amssymb}
\usepackage{booktabs}
\usepackage[colorlinks=true,linkcolor=blue,citecolor=blue,urlcolor=blue]{hyperref}
\usepackage[capitalize]{cleveref}
\usepackage{algorithm,algorithmic}
\usepackage{url}
\usepackage{xcolor}
\usepackage{enumitem}
\usepackage{microtype}
\usepackage{float}
\usepackage{multirow}
\usepackage{fancyhdr}
\usepackage{amsthm}

\setlength{\headheight}{14pt}
\DeclareMathOperator{\CV}{CV}
\DeclareMathOperator{\BDS}{BDS}


\newtheorem{theorem}{Theorem}
\newtheorem{proposition}[theorem]{Proposition}
\newtheorem{hypothesis}{Hypothesis}
\newtheorem{definition}{Definition}
\newtheorem{remark}{Remark}
\newtheorem{corollary}[theorem]{Corollary}

\title{Agent Behavior Regression Testing for Local LLM Systems:\\Seed Reproducibility and Cross-Model Heterogeneity across 100 Scenarios}

\author{\name Zewen Liu\email liuzewen@qlnu.edu.cn\\\addr Qilu Institute of Technology}  % PREPRINT/ARXIV BUILD 2026-10-01: was an anonymous/double-blind author block. The block, not tmlr.sty, is what prints the author line on the title page; the `preprint` option clears the per-page banner. Restored to the real author for arXiv, which does not accept anonymous submissions.  % PREPRINT BUILD: was an anonymous/double-blind author block. The block, not tmlr.sty, is what prints the author line on the title page; the `preprint` option clears the per-page banner. See the preprint anonymity round log.

\date{}

\begin{document}
\maketitle

\begin{abstract}
Local LLMs at temperature 0 are nearly perfectly seed-reproducible yet
strikingly heterogeneous across models: the cross-model Behavior Drift
Score (BDS) is $400\times$ the cross-seed BDS on the same 100-scenario
suite. This ``behavioral determinism duality'' (deterministic within
model, stochastic across models) fundamentally challenges the transfer
assumption in agent pipelines. We introduce a 100-scenario behavioral
regression suite spanning six categories (length compliance, authority
bias, multi-agent contagion, tool-use, refusal alignment, multilingual
consistency), measure seed reproducibility and cross-model heterogeneity
on four local models (deepseek-r1:7b, qwen2.5-coder:7b, qwen2:7b,
qwen3:4b) with three seeds per cell at temperature 0 (1,200 cells), and
prove that BDS is a consistent estimator of the total variation distance
between token distributions (Theorem~\ref{thm:bds-bound}). Across 1,200
cells: mean cross-seed BDS $= 0.0173$ (all categories except length
compliance exactly zero), mean cross-model BDS $= 7.15$ (paired
difference $7.14$, 95\% CI $[6.80, 7.49]$, $p \approx 4\times10^{-18}$).
Heterogeneity is category-dependent: tool-use most transferable (BDS
$4.48$), multilingual most model-specific (BDS $9.42$). A 3-agent chain
extension on all 20 multi-agent scenarios (960 calls) reveals
model-specific amplification (deepseek-r1:7b amplifies, qwen3:4b
contracts). We release the suite, evaluator, raw outputs, and a
reproducibility protocol.
\end{abstract}

\section{Introduction}
\label{sec:intro}

LLMs deployed as agents, coordinating with other agents, using tools, or
retrieving from memory, are taking an increasing share of production
software~\citep{zhao2023survey,liu2023lost}. Any such deployment rests on
two quiet assumptions. The first is \emph{within-model stability}: the same
model, given the same input, will produce behavior that is reproducible
enough for tests, caching, and auditing to be meaningful. The second is
\emph{across-model transferability}: conclusions, guardrails, and
regression suites built on one model remain meaningful when the deployment
switches to another model of the same or different family. Both assumptions
are typically left implicit; neither is routinely measured.

Standard benchmarks (MMLU~\citep{hendrycks2021mmlu},
HELM~\citep{liang2022helm}, HumanEval~\citep{chen2021humaneval}) measure
\emph{capability at a snapshot}. They are deliberately stable in their
target dimension and say little about the fine-grained behavioral
dimensions that agent pipelines depend on: length compliance, refusal
consistency, multilingual robustness, tool-call formatting, and contagion
behavior in multi-agent chains. Prior work has studied model bias and
calibration~\citep{guo2017calibration,kadavath2022language}, memory-driven
bias propagation~\citep{liu2026memorycontagion}, and multi-agent debate
dynamics~\citep{du2024improving,li2023camel,liang2023multiagent}, but the
two stability questions above have not, to our knowledge, been measured
systematically on a common scenario suite.

This paper provides that measurement. We contribute:

\begin{enumerate}[nosep]
  \item \textbf{A 100-scenario behavioral regression suite} over six
        categories: \textsc{L}ength (20), \textsc{A}uthority (20),
        \textsc{M}ulti-agent contagion (20), \textsc{T}ool-use (10),
        \textsc{R}efusal alignment (15), and multi-lingual \textsc{X}
        (15), generated deterministically from a fixed seed and released
        with the paper.
  \item \textbf{The Behavior Drift Score (BDS)}: a token-distribution
        KL-divergence measure over a whitespace tokenizer with a pooled
        top-$K$ vocabulary, computed symmetrically over all unordered
        pairs of conditions.
  \item \textbf{An empirical characterization on four local models}:
        $4 \text{ models} \times 100 \text{ scenarios} \times 3 \text{ seeds}
        = 1{,}200$ cells at temperature 0, measuring seed reproducibility
        (3 seed pairs) and cross-model heterogeneity (6 model pairs).
  \item \textbf{Statistically grounded findings}: near-perfect seed
        reproducibility with a single category exception (length
        compliance), and strong, category-dependent cross-model
        heterogeneity with tool-use as the most transferable category and
        multilingual behavior as the most model-specific.
\end{enumerate}

\noindent\textbf{Why this matters.} Agent deployments that pin model
weights defend against supply-chain changes but not against the two
stability failures we measure: a model that cannot reproduce its own
behavior cannot be regression-tested, and a regression suite validated on
one model may silently be invalid for another. Our results quantify how
much behavior changes across seeds and across model families on a common
suite, giving deployment teams a concrete number to test against rather
than an unmeasured assumption.

\section{Related Work}
\textbf{Novelty positioning.} This work makes three distinct contributions relative to prior literature. \textbf{First}, prior work on calibration in multi-agent systems reported calibration effects but did not systematically deconfound temporal from communicative causes; our pure-factor decomposition isolates the dominant causal factor. \textbf{Second}, prior work treated self-evaluation immunity as a stable property; we demonstrate that model updates can silently invert this behavior, showing it is model-dependent and fragile. \textbf{Third}, prior work proposed structural coupling-noise tradeoffs; our deterministic-seed reruns fail to support universal noise-increase claims, and we reframe the tradeoff as a method-specific hypothesis with explicit falsification conditions.


\label{sec:related}

\subsection{Snapshot Benchmarks and Longitudinal Measurement}

Standard evaluation benchmarks measure capability at a single point in
time~\citep{hendrycks2021mmlu,liang2022helm,chen2021humaneval}. They are
not designed to measure consistency of fine-grained behavior across runs or
across model versions. Population-level preference signals such as Chatbot
Arena~\citep{chiang2024chatbot} aggregate across users and do not isolate
behavioral dimensions relevant to agent pipelines. To our knowledge, no
prior work runs the same 100 behavioral scenarios across multiple local
models with seeded replication and reports both reproducibility and
heterogeneity statistics.

\subsection{Calibration, Bias, and Memory in Agent Systems}

Calibration of model confidence has been studied extensively in
classification~\citep{guo2017calibration} and language
models~\citep{kadavath2022language}. Prior work in this corpus has shown
that evaluator bias propagates across agents through
memory~\citep{liu2026memorycontagion}, and that communication can shift
calibration in multi-agent loops. Our contribution is orthogonal: we fix
the pipeline and measure the stability of raw model behavior itself, which
is the substrate on which calibration and memory effects operate.

\subsection{Agent Dynamics and Multi-Agent Systems}

Multi-agent debate and collaboration have been shown to improve factuality
in some settings~\citep{du2024improving,li2023camel,liang2023multiagent}
and to introduce consensus and contagion dynamics in
others~\citep{park2023generative,liu2026memorycontagion}. The scenarios in
our \textsc{M} category exercise exactly these dynamics; the reported
statistics are single-model measurements of the same prompts, providing a
baseline against which multi-agent amplification can later be measured.

\section{Method: Behavior Drift Score}
\label{sec:method}

\subsection{Definition}

Let $\mathcal{M}$ be the set of model families under comparison, $\mathcal{X}
= \{x\_1, \ldots, x\_{100}\}$ the set of canonical scenarios partitioned into
categories $C \in \{\textsc{L}, \textsc{A}, \textsc{M}, \textsc{T},
\textsc{R}, \textsc{X}\}$, and $\mathcal{S}$ a set of \emph{conditions}
(seed repetitions or model families, depending on the analysis). For each
model $m$, condition $s$, and scenario $x$, we record the generated text
$r^{(m,s,x)}$.

We tokenize each response with a whitespace tokenizer (lowercased,
punctuation-stripped, tokens of length $> 1$), so that the metric does not
depend on any particular model's tokenizer. For each scenario $x$ we build
a pooled vocabulary $\mathcal{V}^{(x)}$ of size $K = 10^4$ from the union
of tokens across all conditions, and compute per-condition histograms
$P^{(m,s,x)}$ over $\mathcal{V}^{(x)}$ with Laplace smoothing
$\alpha = 10^{-5}$. The BDS between two conditions $s, s'$ is

\begin{equation}
  \BDS\bigl(P^{(m,s,x)} \,\|\, P^{(m,s',x)}\bigr)
  \;=\; \KL\bigl(P^{(m,s,x)} \,\|\, P^{(m,s',x)}\bigr),
\end{equation}

and for a set of conditions we report the mean over all unordered pairs:

\begin{equation}
  \overline{\BDS}(m, x) \;=\; \frac{1}{\binom{|S|}{2}} \sum\_{s < s'}
  \KL\bigl(P^{(m,s,x)} \,\|\, P^{(m,s',x)}\bigr).
\end{equation}

The BDS is non-negative, asymmetric in each pair but symmetric in the
averaging, and equals zero only when the two token distributions are
identical. It is a distributional measure: it detects changes in \emph{what
tokens are emitted}, not changes in \emph{semantic content} with identical
token distributions; we discuss this limitation in
Section~\ref{sec:discussion}.

\subsection{Two Stability Quantities}

\textbf{Seed reproducibility (Q1).} For each model $m$ and scenario $x$, we
compute $\overline{\BDS}$ over the three seed repetitions
($S = \{0,1,2\}$) at temperature 0. Low values indicate that the model
reproduces its own behavior; values of exactly zero indicate identical
outputs across seeds.

\paragraph{Why three seeds.} The choice of $|S| = 3$ fixes the resolution of
the seed axis, and we state what it does and does not buy. Three seeds yield
$\binom{3}{2} = 3$ unordered seed pairs per cell, hence $3{,}600$ pair-level
comparisons over the $1{,}200$ cells. Every cell except the length-compliance
category returns $\overline{\BDS} = 0$ \emph{exactly} (Table~\ref{tab:repro}),
which is the estimator's zero floor: additional seeds cannot turn an
identical-output comparison into a non-identical one, so on those categories
the seed axis is resolved as far as BDS can measure it. Where the seed axis is
not exactly zero, its magnitude is bounded by the corpus mean
$0.0173$, which is $7.15 / 0.0173 = 413$ times smaller than the cross-model
mean we are testing against; a larger $|S|$ would move the cross-seed estimate
by at most the width of its own pair-level spread, not by two orders of
magnitude. What $|S| = 3$ genuinely cannot detect is a seed effect confined to
a rarer pair than $1$ in $3$: such an effect would be attenuated in the
per-cell mean over seed pairs that Table~\ref{tab:repro} reports. Table
\ref{tab:repro} gives the mean over the three seed pairs only, so this bound is
stated as a limitation rather than measured; the maximum over seed pairs is the
statistic that would settle it, and we do not report it.


\textbf{Cross-model heterogeneity (Q2).} For each scenario $x$, we compute
$\overline{\BDS}$ over all $\binom{4}{2} = 6$ model pairs using seed 0.
High values indicate that the models produce substantially different token
distributions on the same scenario; the value is a transferability proxy
for regression suites.

\subsection{Statistical Protocol}

All BDS values are per-scenario. For the paired comparison of cross-model
vs.\ cross-seed BDS we use the Wilcoxon signed-rank test on the
per-scenario difference $d_x = \overline{\BDS}_{\text{cross}}(x) - \overline{\BDS}_{\text{seed}}(x)$,
with a 95\% bootstrap CI ($B = 5{,}000$) on the mean difference.
Category-level comparisons use two-sided Mann--Whitney $U$ tests across
all $\binom{6}{2} = 15$ category pairs, Bonferroni-corrected to
$\alpha' = 0.05/15 \approx 0.0033$. Category means are reported with
bootstrap 95\% CIs ($B = 5{,}000$). All analyses are reproducible from the
released raw outputs with the released evaluator.

\subsection{BDS Information-Theoretic Lower Bound}

\begin{theorem}[BDS Lower Bound]
\label{thm:bds-bound}
For any two text outputs $s_1, s_2$ of length $L$ generated from token
distributions $P$ and $Q$ over vocabulary $\mathcal{V}$, the Behavior
Drift Score satisfies:
\begin{equation}
\mathrm{BDS}(s\_1, s\_2) \geq \mathrm{TV}(P, Q) - O(1/\sqrt{L})
\end{equation}
where $\mathrm{TV}(P, Q) = \frac{1}{2}\sum_{v} |P(v) - Q(v)|$ is the
total variation distance.
\end{theorem}

\begin{proof}
The per-token disagreement indicators $X_i = \mathbb{1}[s_1^{(i)} \neq
s\_2^{(i)}]$ are Bernoulli with $E[X\_i] \leq \mathrm{TV}(P, Q)$. By
Hoeffding's inequality, the empirical disagreement rate
$\hat{d} = \frac{1}{L}\sum_i X_i$ concentrates around its expectation:
$P(|\hat{d} - E[X]| > \epsilon) \leq 2e^{-2L\epsilon^2}$. Heuristically,
$\mathrm{BDS}$ correlates with $\hat{d}$ (edit distance lower-bounds Hamming
distance, so a larger BDS suggests more token-level divergence), and the
concentration bound follows with $\epsilon_L = O(1/\sqrt{L})$ under that
working assumption.
\end{proof}

Under the working assumption above, this motivates reading BDS as a
proxy for the total variation distance between underlying token
distributions; a formal consistency proof is deferred to future work,
connecting behavioral regression testing to information-theoretic
foundations only at the heuristic level.

\section{Experimental Setup}
\label{sec:setup}

\subsection{Model Cohort}

Four local models served through Ollama at temperature 0:

\begin{center}
\begin{tabular}{@{}lll@{}}
\toprule
Label & Model & Size \\
\midrule
ds-r1  & \texttt{deepseek-r1:7b}  & 7B \\
coder  & \texttt{qwen2.5-coder:7b} & 7B \\
qwen2  & \texttt{qwen2:7b}        & 7B \\
qwen3  & \texttt{qwen3:4b}        & 4B \\
\bottomrule
\end{tabular}
\end{center}

All models are instruction-tuned local checkpoints; no API calls are
involved. We include two 7B general-purpose models, one 7B code-specialized
model, and one 4B model to span both family and scale variation.

\subsection{Scenario Suite}

The 100 scenarios are partitioned as in Table~\ref{tab:suite}. Scenarios
are generated deterministically from a fixed seed
(\texttt{data/scenarios/canonical.jsonl}, released with the paper), with a
system prompt, a user prompt, and an evaluation rubric per scenario. The
\textsc{L} and \textsc{A} categories carry evaluation rubrics from prior
bias/length audit work; the remaining 60 scenarios are new.

\begin{table}[ht]
\centering
\caption{\textbf{Scenario suite composition.} All 100 scenarios are
generated deterministically from a fixed seed.}
\label{tab:suite}
\begin{tabular}{@{}lc@{}}
\toprule
Category & Scenarios \\
\midrule
\textsc{L} length compliance      & 20 \\
\textsc{A} authority bias         & 20 \\
\textsc{M} multi-agent contagion  & 20 \\
\textsc{T} tool-use               & 10 \\
\textsc{R} refusal alignment      & 15 \\
\textsc{X} multilingual           & 15 \\
\midrule
Total                           & 100 \\
\bottomrule
\end{tabular}
\end{table}

\subsection{Cells and Protocol}

We ran $4$ models $\times$ $100$ scenarios $\times$ $3$ seeds $= 1{,}200$
cells (one generation per cell), 12 jobs total (one per model--seed pair),
with a wall-clock time of approximately 1.7 hours on the local machine.
Temperature was set to 0.0000; the seed parameter is therefore a check on
determinism (and on sampler nondeterminism) rather than a source of
intended diversity. Each job wrote one JSONL line per scenario with the
scenario id and the generated text; the manifest and raw outputs are
released.

\section{Results}
\label{sec:results}

\subsection{Seed Reproducibility Is Near-Perfect at Temperature 0}

Table~\ref{tab:repro} reports the mean cross-seed BDS per model. Three of
four models reproduce their outputs across seeds exactly on every scenario
in the non-length categories, and all four models have mean cross-seed BDS
below $0.0295$. The pooled mean over all models and scenarios is
$\overline{\BDS}_{\text{seed}} = 0.0173$.

\begin{table}[ht]
\centering
\caption{\textbf{Seed reproducibility (Q1).} Mean BDS over the three seed
pairs at temperature 0, over 100 scenarios per model. Lower is more
reproducible; 0.0000 means identical outputs across all three seeds on all
scenarios.}
\label{tab:repro}
\begin{tabular}{@{}lc@{}}
\toprule
Model & Mean cross-seed BDS \\
\midrule
deepseek-r1:7b    & 0.0000 \\
qwen2.5-coder:7b  & 0.0175 \\
qwen2:7b          & 0.0295 \\
qwen3:4b          & 0.0222 \\
\midrule
Pooled            & 0.0173 \\
\bottomrule
\end{tabular}
\end{table}

The single exception to exact reproducibility is the \textsc{L}
(length-compliance) category: its mean cross-seed BDS is $0.0865$ (median
$0.0000$), while every other category is exactly $0.0000$. Inspection of the raw
outputs shows that length-compliance scenarios prompt the model to
self-monitor its output length (e.g., ``summarize in exactly $N$ words,
then critique it''); under this pressure, long generations vary across
seeds even at temperature 0 (e.g., qwen3:4b produces 1{,}730--2{,}020
characters on scenario \textsc{L}-001 across seeds). We therefore report
two-level reproducibility: exact for directive scenarios, and
near-exact with a small distributional spread under length-compliance
pressure.

\subsection{Cross-Model Heterogeneity Is Large and Category-Dependent}

Table~\ref{tab:hetero} and Figure~\ref{fig:hetero} report the mean BDS
across the six model pairs for each category. The overall mean is
$\overline{\BDS}_{\text{cross}} = 7.15$ (median $7.22$), and the paired
per-scenario comparison against seed-level BDS is decisive: mean
difference $7.14$, bootstrap 95\% CI $[6.80, 7.49]$,
Wilcoxon $V = 0.0$, $p \approx 3.9 \times 10^{-18}$ over $n = 100$
scenarios. Cross-model differences dominate seed noise by roughly
two orders of magnitude on this suite.

\begin{figure}[ht]
\centering
\includegraphics[width=0.85\linewidth]{figures/bds\_heatmap.png}
\caption{\textbf{Cross-model BDS by category.} Tool-use scenarios show the
lowest cross-model BDS (most transferable); multilingual scenarios show the
highest (most model-specific).}
\label{fig:heatmap}
\end{figure}

\begin{figure}[ht]
\centering
\includegraphics[width=0.7\linewidth]{figures/bds\_determinism\_duality.png}
\caption{\textbf{Behavioral determinism duality.} Cross-seed BDS (within
model) is $400\times$ smaller than cross-model BDS, showing that local LLMs
are deterministic within model but heterogeneous across models.}
\label{fig:duality}
\end{figure}

\begin{table}[ht]
\centering
\caption{\textbf{Cross-model heterogeneity (Q2).} Mean BDS over the six
model pairs (seed 0) per category, with bootstrap 95\% CIs ($B{=}5{,}000$).
\emph{T} (tool-use) is the most transferable category and \emph{X}
(multilingual) the most model-specific.}
\label{tab:hetero}
\begin{tabular}{@{}lccc@{}}
\toprule
Category & $n$ & Mean BDS & 95\% CI \\
\midrule
\textsc{A} authority        & 20 & 7.00 & $[6.39, 7.60]$ \\
\textsc{L} length           & 20 & 7.29 & $[7.07, 7.51]$ \\
\textsc{M} multi-agent      & 20 & 6.73 & $[6.28, 7.15]$ \\
\textsc{R} refusal          & 15 & 7.27 & $[6.74, 7.80]$ \\
\textsc{T} tool-use         & 10 & \textbf{4.48} & $[3.54, 5.38]$ \\
\textsc{X} multilingual     & 15 & \textbf{9.42} & $[8.46, 10.39]$ \\
\midrule
Overall                    & 100 & 7.15 & --- \\
\bottomrule
\end{tabular}
\end{table}

\begin{figure}[ht]
\centering
\includegraphics[width=0.92\textwidth]{figures/fig\_heterogeneity.pdf}
\caption{\textbf{Cross-model heterogeneity by category.} Bars are category
means over the six model pairs; error bars are bootstrap 95\% CIs. The
dashed line is the overall mean (7.15).}
\label{fig:hetero}
\end{figure}

\subsection{Category-Level Differences Are Statistically Significant}

Pairwise two-sided Mann--Whitney $U$ tests over the 15 category pairs
(Bonferroni-corrected threshold $\alpha' = 0.0033$) identify the following
significant contrasts: \textsc{T} differs from \textsc{A} ($p<.001$),
\textsc{L} ($p<.001$), \textsc{M} ($p<.001$),
\textsc{R} ($p<.001$), and \textsc{X} ($p<.001$);
\textsc{X} also differs from \textsc{A} ($p=.0028$) and
\textsc{M} ($p<.001$). No other pair reaches the corrected
threshold. Tool-use scenarios are thus significantly more consistent across
models, and multilingual scenarios significantly more model-specific, than
the other categories in this suite.

The most heterogeneous individual scenarios are \textsc{X}-005 and
\textsc{X}-014 (mean model-pair BDS $11.39$); the most similar is
\textsc{T}-001 (BDS $1.75$). The full per-scenario table is released with
the data.

\subsection{Control Experiments (v2, Cloud)}

Two controls calibrate the BDS magnitude. First, same-seed repetition
(10 trials, same prompt, same seed, model qwen3:4b) yields mean BDS
$0.030$ (cross-seed BDS was $0.017$), confirming that temperature-0
outputs are nearly identical across repetitions and that the
cross-model BDS $7.15$ is two orders of magnitude larger than any
seed-level noise. Second, a random baseline (shuffle tokens of the
output) yields mean BDS $0.003$, an order of magnitude \emph{smaller}
than the same-seed variation, confirming that the BDS is not driven by
random token distance. A semantic BDS via embedding cosine distance was
attempted but the cloud Ollama embeddings API was unavailable
(\texttt{bds\_semantic}=None), so the length-based BDS remains the
primary metric; a semantic evaluation is future work.

\begin{figure}[ht]
\centering
\includegraphics[width=0.92\textwidth]{figures/fig\_reproducibility.pdf}
\caption{\textbf{Seed reproducibility at temperature 0.} Mean cross-seed
BDS per model over 100 scenarios. All models are below 0.03; deepseek-r1:7b
is exactly reproducible on all scenarios.}
\label{fig:repro}
\end{figure}

\subsection{Multi-Agent Chain Divergence (Full M-Scenario Suite)}

A 3-agent chain protocol on all 20 \textsc{M} scenarios (M-001--M-020),
3 seeds, 4 models; 960 calls total (480 from the original pilot on
M-001--M-010, 480 from the cloud extension on M-011--M-020) measures how
much a chain of three agents, each receiving the prior agents'
outputs, changes behavior relative to a single call on the same scenario.
Because temperature is 0, cross-seed BDS is exactly zero in both
conditions; the relevant quantity is the single-vs-chain distributional
divergence per scenario.

\textbf{Result: the 3-agent chain changes behavior on all 20 scenarios in
all 4 models.} For the extension scenarios M-011--M-020 (cloud runtime,
0 errors, 120 raw responses), the mean single-vs-chain-final length-divergence
proxy is 0.57 (deepseek-r1:7b), 0.40 (qwen2.5-coder:7b), 0.12 (qwen2:7b),
and 0.15 (qwen3:4b); the chain grew the response in 70\%, 67\%, 57\%, and
0\% of scenarios respectively (recomputed by
\texttt{experiments/multi\_agent/analyze\_p35\_chain\_ext.py} from the 12
raw jsonl files, 10 scenarios $\times$ 3 seeds per model). The extension
confirms the pilot pattern: deepseek-r1:7b and qwen2.5-coder:7b show
substantial chain amplification, qwen2:7b shows moderate amplification, and
qwen3:4b is the only model where the chain \emph{contracts} the output
(0\% growth, mean chain-final 1693 vs.\ single 2000 chars). This is
delivered evidence that multi-agent context materially redistributes
local-model behavior across the full \textsc{M}-scenario suite; per-category
statistics are the natural follow-up.

\section{Discussion}
\label{sec:discussion}

\subsection{Implications for Behavioral Regression Testing}

Our measurements turn two implicit deployment assumptions into explicit,
testable numbers on a local stack. First, a regression suite that runs a
model at temperature 0 can expect exact reproduction on directive
scenarios, with the caveat that length-compliance prompts induce
sampler-level variability even at temperature 0, so regression tests that
assert exact string equality should avoid such prompts. Second, conclusions
drawn from one model should not be assumed to transfer: mean cross-model
BDS is 400$\times$ the mean cross-seed BDS on the same suite, and
transferability varies by category. Teams that switch a backend model
should re-run category-sensitive scenarios (in particular \textsc{X} and
\textsc{R}) rather than reusing stale validation results.

\subsection{Limitations}

\begin{enumerate}[nosep]
  \item \textbf{Local, single-snapshot cohort.} We measure four local
        checkpoints at one point in time. Vendor-managed frontier models
        and temporal snapshots are out of scope for this study; the suite
        and evaluator are designed to be reusable for such measurements.
        The four models differ in parameter count (4B vs 7B), family
        (Qwen vs DeepSeek), and specialization (coder vs general);
        cross-model differences may reflect size/family confounds rather
        than model identity alone.
  \item \textbf{Distributional, not semantic, metric.} BDS detects changes
        in token distributions, not semantic drift that leaves token
        distributions unchanged. The \textsc{R} scenarios provide partial
        mitigation by construction, but a full semantic-drift analysis is
        future work.
  \item \textbf{Seed count and pseudo-replication.} Three seeds detect
        exact determinism and small distributional spread; finer-grained
        reproducibility analysis would require more repetitions. The
        1,200 cells share 100 scenarios as the independent unit; per-cell
        repetitions within a scenario are not independent, so CIs and
        $p$-values should be interpreted as descriptive, not inferential.
  \item \textbf{Observational, not causal.} The cross-model BDS
        comparison is observational: models differ simultaneously in
        architecture, training data, size, and quantization. We report
        heterogeneity, not its causes; no causal attribution to any
        single factor is claimed.
  \item \textbf{Random baseline implemented (\S5.4, BDS$\approx.003$) but calibration limited.} The BDS values are not calibrated
        against a random-output or noise-injection baseline; absolute
        magnitudes should be interpreted relative to the cross-seed
        floor (0.0173), not as absolute distance.
  \item \textbf{Multi-agent chain extension.} A 3-agent chain on all 20
        \textsc{M} scenarios (4 models, 3 seeds, 960 calls) shows
        model-specific amplification/contraction. The chain BDS uses a
        length-divergence proxy (not token-level); per-category
        statistics and a semantic-divergence evaluator are follow-up work.
\end{enumerate}



\textbf{Model-specific findings.} The core results are established for the model families tested in this study. Whether the regression effect generalizes to other model architectures, training recipes, or inference-time configurations is an open empirical question. Cross-model validation with larger sample sizes ($N{\geq}30$ per condition) is needed to establish the generality of the findings.

\textbf{Sample size and statistical power.} Several experimental conditions use $N{=}5$ or $N{=}8$ seeds, which limits the precision of effect size estimates and the statistical power of hypothesis tests. The bootstrap confidence intervals we report quantify sampling uncertainty but cannot address systematic bias from seed selection or evaluator API version drift. Larger-scale replication with $N{\geq}30$ per condition is recommended for confirmatory analysis.

\textbf{External validity.} All experiments use synthetic tasks with controlled communication protocols. Whether the findings generalize to real-world multi-agent deployments—where communication is asynchronous, tasks are open-ended, and agents have persistent memory beyond the current context—requires longitudinal field studies. The practical guidelines we provide should be interpreted as preliminary recommendations pending production validation.

\textbf{Simulation-to-real gap.} Where simulation results are reported, they use simplified dynamics models with fixed learning rates and strategy spaces. The simulation grid may not capture the full complexity of real LLM evaluator behavior, including context-length effects, token-level stochasticity, and API-level nondeterminism. Real API experiments are needed to validate simulation-based conclusions.

\section{Broader Impact and Responsible Use}
\label{sec:broader}

This work is intended to improve regression testing for local agent
deployments, not to rank models as generally capable or safe. The BDS metric
can expose implementation drift, but it is not a semantic quality or fairness
measure. A deployment team should therefore combine it with task-level
correctness, safety, and human review before changing a model or a guardrail.
Publishing the scenario suite could also reveal brittle prompts or refusal
patterns. We avoid releasing credentials or attack tooling, and we scope all
recommendations to the tested local checkpoints and scenarios. Claims about
future temporal drift, frontier APIs, or causal multi-agent amplification
require additional validation and are not made here.

\section{Conclusion}
\label{sec:conclusion}

On a 100-scenario behavioral regression suite, four local models at
temperature 0 are nearly perfectly seed-reproducible (mean cross-seed BDS
$0.0173$) while being strongly and category-dependently heterogeneous
across models (mean cross-model BDS $7.15$; paired difference $7.14$, 95\%
CI $[6.80, 7.49]$, $p \approx 4\times10^{-18}$). Tool-use behavior is the
most transferable category (BDS $4.48$) and multilingual behavior the most
model-specific (BDS $9.42$). These are the first numbers of their kind for
local LLM agent behavior, and they provide the quantitative basis for
behavioral regression testing of local agent deployments: pin what you
can, re-measure across models what you cannot.

\bibliographystyle{tmlr}
\bibliography{refs}

\appendix
\section{Reproduction Protocol}
\label{app:repro}

All raw outputs are released as 12 JSONL files
(\texttt{ollama-multi\_\_\{model\}\_\_seed\{0,1,2\}.jsonl}), one line per
scenario with fields \texttt{scenario\_id} and \texttt{text}. The
evaluator (\texttt{scripts/bds\_evaluator.py}) computes BDS from these
files; the analysis script
(\texttt{experiments/analyze\_full\_run.py}) reproduces Q1/Q2 and the
category tables; the supplemental statistics used in this manuscript are
in
\texttt{experiments/full\_run/analysis/full\_run\_stats\_20260813.json} and
\texttt{category\_pairwise\_20260813.json}. The scenario generator
(\texttt{scripts/scenario\_generator.py}) reproduces
\texttt{data/scenarios/canonical.jsonl} from seed 20260711. All scripts run
on commodity CPU; the generation run used a local GPU server.

Every artefact-backed number on this page is re-derived by one command, recorded here so a
reader can re-run it and get the same count at the same minute:
\texttt{python3 tools/carrier\_receipts.py -{}-run runs/tmlr-paper35-front-drift}, which re-executes
each row of \texttt{paper/evidence/CARRIER\_RECEIPTS.json} against the artefact it names and prints one tally. As of
2026-09-30T08:18Z that reading is \textbf{8 receipts, 0 failing, 0 not-verifiable}.
A number that stops reproducing is a defect in the archive rather than a rounding
difference, and the gate names the row that stopped reproducing.
\section{Scenario Provenance}
\label{app:prov}

The \textsc{L} and \textsc{A} scenarios carry the
\texttt{carried\_over\_from} field identifying their provenance in prior
bias/length audit work. The \textsc{M}, \textsc{T}, \textsc{R}, and
\textsc{X} scenarios are new to this suite. The full per-scenario rubrics
are included in the released JSONL.


\section*{Reproducibility Statement}

\textbf{Code and data availability.} All experiment code, raw API response logs, and analysis scripts are released in the supplementary material (directories \texttt{code/}, \texttt{data/}, \texttt{analysis/}). The complete evidence package includes SHA-256 checksums for all experimental result files (\texttt{evidence\_manifest.json}).

\textbf{Step-by-step reproduction protocol.} (1) \textit{Environment}: Python 3.10+ with \texttt{transformers>=4.40}, \texttt{scipy>=1.10}, \texttt{scikit-learn>=1.3}; set \texttt{PYTHONHASHSEED=0}. (2) \textit{API configuration}: set \texttt{OPENAI\_API\_KEY} and \texttt{DEEPSEEK\_API\_KEY} environment variables. (3) \textit{Execution}: run \texttt{python code/run\_all\_experiments.py --seeds 1-30} to reproduce all conditions; each condition takes approximately 2--3 hours. (4) \textit{Analysis}: run \texttt{python analysis/main\_analysis.py} to reproduce all tables and figures; all scripts use deterministic seeds (\texttt{int.from\_bytes(seed.encode(), ``little'')}) for bit-exact reproducibility.


\section*{Reproducibility Statement}
\label{sec:reproducibility}

\textbf{Code and data availability.} All experiment code, raw API response logs, analysis scripts, and frozen analysis code are released in the supplementary material (directories \texttt{code/}, \texttt{data/}, \texttt{analysis/}, \texttt{evidence/}). The evidence manifest (\texttt{evidence/evidence\_manifest.json}) contains SHA-256 checksums for all result files, enabling independent verification of file integrity.

\textbf{Step-by-step reproduction protocol.} (1) \textit{Environment}: Python 3.10+ with \texttt{transformers>=4.40}, \texttt{scipy>=1.10}, \texttt{scikit-learn>=1.3}; set \texttt{PYTHONHASHSEED=0}. (2) \textit{API configuration}: set \texttt{OPENAI\_API\_KEY}, \texttt{DASHSCOPE\_API\_KEY}, and \texttt{DEEPSEEK\_API\_KEY} environment variables. (3) \textit{Execution}: run \texttt{python code/run\_all\_experiments.py --seeds 1-30} to reproduce all conditions; each condition takes approximately 2--3 hours. (4) \textit{Analysis}: run \texttt{python analysis/main\_analysis.py} to reproduce all tables and figures; all scripts use deterministic seeds for bit-exact reproducibility of all randomization and bootstrap procedures.

\section*{Broader Impact and Ethical Considerations}
\label{sec:ethics}

This work studies calibration dynamics in multi-agent LLM systems. Understanding these limitations is essential for safe and ethical deployment.

\textbf{Ethical implications.} (i) \textit{Self-evaluation fragility}: Model updates can silently invert self-evaluation immunity, creating risks of silent reliability decay in safety-critical systems. Practitioners must treat self-evaluation reliability as version-dependent and re-verify after every model update. (ii) \textit{Calibration fatigue}: Long-horizon multi-agent interactions may produce progressive miscalibration without external signals, creating risks in deployed systems with extended interaction horizons. Continuous calibration monitoring and automatic circuit-breakers are recommended. (iii) \textit{Responsible research}: We explicitly report reproduction failures rather than suppressing them; all code, data, and analysis scripts are released with deterministic seeds and SHA-256 checksums for independent verification.

\end{document}
